\documentclass[11pt,a4paper]{article}
\usepackage{amsmath,amssymb}
\DeclareMathOperator{\Tr}{Tr}
\DeclareMathOperator{\Var}{Var}
\DeclareMathOperator{\Cov}{Cov}
\begin{document}
\title{Observational entropy, QFI, and a possible Pesin connection}
\author{}
\date{}
\maketitle
\section*{1. The exact connection}
Assume a finite-dimensional system and a fixed POVM $E_a$, with $p_a=\Tr(\rho E_a)>0$, $v_a=\Tr E_a$, and natural logarithms. Introduce the score $u_a=\dot p_a/p_a$ and $\ell_a=\log(p_a/v_a)$. Normalization implies $\sum_ap_au_a=0$, so
\[S_E=-\sum_ap_a\ell_a,\qquad\dot S_E=-\Cov_p(u,\ell).\]
Classical Fisher information is $I_E^{(t)}=\sum_ap_au_a^2$. Cauchy--Schwarz and measurement monotonicity of SLD quantum Fisher information give
\[\boxed{|\dot S_E|\le\sqrt{I_E^{(t)}\Var_p(\ell)}\le\sqrt{F_Q^{(t)}\Var_p(\ell)}.}\]
Here the SLD satisfies $\dot\rho=(\rho L+L\rho)/2$ and $F_Q^{(t)}=\Tr(\rho L^2)$. The measurement bound follows by applying Hilbert--Schmidt Cauchy--Schwarz to $\sqrt{E_a}\sqrt\rho$ and $\sqrt{E_a}L\sqrt\rho$, then summing over $a$.

\textbf{Why the connection exists:} entropy rate is the pairing of a probability velocity with an entropy gradient; Fisher information is the squared statistical speed. Speed alone does not determine the sign or magnitude of the entropy derivative. Equality additionally requires alignment of the score with centered $\ell_a$ and a measurement attaining QFI.

The argument needs no density-matrix coordinates. For arbitrary state coordinates $x^\mu$ and an observation model $p_a(x)$,
\[g^E_{\mu\nu}=\sum_a\frac{\partial_\mu p_a\,\partial_\nu p_a}{p_a},\qquad I_E^{(t)}=g^E_{\mu\nu}\dot x^\mu\dot x^\nu.\]
This observation metric can be degenerate. Zero probabilities require limiting arguments; smooth unitary pure-state trajectories have quadratic opening of zero outcomes, sufficient for the entropy-rate formula. Mere differentiability of an arbitrary density curve is not sufficient.

\section*{2. Why pure-state optimization gives energy uncertainty}
For a normalized pure state, remove the unphysical phase velocity:
\[|w\rangle=(I-|\psi\rangle\langle\psi|)|\dot\psi\rangle,\qquad s=\|w\|.\]
Taking $|0\rangle=|\psi\rangle$ and $|1\rangle=|w\rangle/s$ for $s>0$ gives
\[\dot\rho=s(|1\rangle\langle0|+|0\rangle\langle1|),\qquad F_Q^{(t)}=4s^2.\]
All pure state--tangent pairs with equal $s$ are unitarily equivalent. Conjugating a POVM preserves its probabilities, tangent probabilities, and trace volumes in the corresponding transformed pair. Since the rate is linear in $s$, optimization over all $m$-outcome POVMs yields
\[\boxed{\sup_E\dot S_E=C_{d,m}s=\frac{C_{d,m}}2\sqrt{F_Q^{(t)}}=C_{d,m}\frac{\Delta H}{\hbar}.}\]
The last equality uses Schr\"odinger evolution; the geometric argument applies to any differentiable pure-state curve. The coefficient is a canonical optimization constant, not a chaos parameter. Its exact numerical value is not established by previous numerical optimizations.

The POVM is held fixed during differentiation, but may be chosen anew for each state and instant. Consequently this is not a sustained entropy-growth law. For independently optimized ensemble members the mean rate is $(C_{d,m}/2)\overline{\sqrt{F_Q}}$, not $(C_{d,m}/2)\sqrt{\overline{F_Q}}$. A single common POVM has only the corresponding upper bound. Mixed states do not generally obey equality with energy uncertainty: a state commuting with $H$ has zero time QFI but may have nonzero energy variance.

\section*{3. The two dynamical equations concern different directions}
For an autonomous classical Hamiltonian in canonical coordinates,
\[\dot z=J\nabla H(z),\qquad\dot\xi=JD^2H(z(t))\xi.\]
The second equation is exact for an infinitesimal initial-condition variation $\xi=D\Phi_t\,\xi_0$. Lyapunov exponents measure growth of such variations. Time QFI instead measures the speed along the orbit. Even classically these are different: $X_H(z(t))=D\Phi_tX_H(z_0)$ is one special tangent direction. On a compact invariant region without equilibria its exponent is zero, although transverse exponents can be positive.

Quantum variables can be wavefunction amplitudes rather than density matrices. Write $H=A+iB$, with $A^{\mathsf T}=A$, $B^{\mathsf T}=-B$, and $\psi=(q+ip)/\sqrt2$. Then
\[\dot x=Kx,\qquad\dot{\delta x}=K\delta x,\qquad K=\frac1\hbar\begin{pmatrix}B&A\\-A&B\end{pmatrix},\quad x=\begin{pmatrix}q\\p\end{pmatrix}.\]
This is a Hamiltonian system, but $K$ is also antisymmetric: its propagator is orthogonal. On normalized rays, unitary evolution preserves the Fubini--Study metric. Thus an initial preparation parameter evolved by a common parameter-independent unitary has exactly conserved QFI:
\[|\psi_t(\theta)\rangle=U_t|\psi_0(\theta)\rangle\quad\Longrightarrow\quad F_Q^{(\theta)}(t)=F_Q^{(\theta)}(0).\]
For autonomous $H$, physical-time QFI is likewise constant, $F_Q^{(t)}=4(\Delta H)^2/\hbar^2$. A periodic qubit with nonzero $\Delta H$ already has positive optimized instantaneous OE rate without chaos. Neither quantity is automatically a Lyapunov exponent.

\section*{4. Where a Pesin-type connection could enter}
\textbf{Exact illustrative model, not a quantum theorem.} Take the classical Hamiltonian $H=\sum_i\lambda_iq_ip_i$, with $\lambda_i>0$, an initial Gaussian covariance $\sigma^2I$, and fixed observation blur $\epsilon^2I$. If $M_t$ is the tangent map, the observed covariance is $C_t=\sigma^2M_tM_t^{\mathsf T}+\epsilon^2I$. Its differential entropy satisfies
\[\dot h_\epsilon=\sum_j a_j(t),\qquad a_j(t)=\frac{\gamma_j\sigma^2e^{2\gamma_jt}}{\sigma^2e^{2\gamma_jt}+\epsilon^2},\]
where $\gamma_j$ runs over all $\lambda_i$ and $-\lambda_i$. Expanding directions remain visible while contraction falls below resolution. Hence, for fixed $\epsilon>0$,
\[\dot h_\epsilon\longrightarrow\sum_i\lambda_i,\qquad I_\epsilon^{(t)}=2\sum_ja_j^2\longrightarrow2\sum_i\lambda_i^2.\]
Without blur the entropy rate is zero. The positive-exponent sum therefore arises from unresolved contraction, not fine-grained volume growth. Moreover scalar Fisher information gives a quadratic sum, not the linear sum needed by Pesin. This unbounded Gaussian example has no normalized invariant Liouville probability on the whole space and is not itself Pesin's theorem.

\textbf{A concrete literature lead.} Goldfriend and Kurchan study entropy of a weak-noise-averaged state, obtaining a response-Gram determinant, not fixed-POVM OE: Phys.\ Rev.\ Research \textbf{3}, 023234 (2021), doi:10.1103/PhysRevResearch.3.023234, printed pp.~4--8, Eqs.~(18)--(29), (42)--(49). For the local perturbation family $\exp[-i\eta^\alpha X_\alpha(t)/\hbar]|\psi_0\rangle$, independently applying the pure-state QFI formula gives
\[F_{\alpha\beta}=\frac2{\hbar^2}\langle\{X_{\alpha,c}(t),X_{\beta,c}(t)\}\rangle.\]
Their endpoint covariance therefore has a multiparameter perturbation-QFI interpretation. A half-log determinant of growing eigenvalues can produce a positive-exponent sum. This is not scalar physical-time QFI; the paper's lower-bound saturation and semiclassical/noise limits require additional assumptions.

\textbf{Remaining gap.} A quantum OE/Pesin claim needs a fixed semiclassical measurement, controlled propagation and entropy errors, specified resolution, an appropriate invariant-measure ensemble, and an intermediate window before saturation. Classical Pesin's formula concerns partition-history KS entropy under smoothness and invariant-measure hypotheses, not an arbitrary final-time measurement entropy. At fixed finite $d$, $S_E\le\log d$, so its asymptotic net rate is zero.

\textbf{Conclusion.} The OE--QFI relation is an exact speed/gradient statement. A possible Pesin relation instead concerns resolved unstable-volume growth. A multiparameter response metric plus a justified coarse-graining limit is the promising bridge; identifying it with our fixed-POVM OE rate remains unproved.
\end{document}
